% Irodov Part V §5.3 originals 143--150; printed pp.211--212.
% Original book answers pp.331--332.
\diffractionproblem{143}{211}
光谱仪的三棱镜在最小偏折状态，底边宽度为 $b$（图5.28），
材料折射率 $n=n(\lambda)$。
利用夫琅禾费衍射与瑞利分辨准则证明极限分辨本领
$\lambda/\delta\lambda=b|\mathrm dn/\mathrm d\lambda|$。
\exsolution
考虑来自棱镜左右极端的两条光线，
它们的光程差在最小偏折方向满足同一波前条件。
波长改变 $\delta\lambda$ 引起折射率改变量
$\delta n=(\mathrm dn/\mathrm d\lambda)\delta\lambda$，
于是两极端出射波新增相位对应光程差
$\delta\Delta=b\,\delta n$。
根据瑞利准则，当一条谱线的中心极大与另一谱线的第一个衍射极小重合时，
两极端新增的光程差约等于一个波长：
$|b\,\delta n|=\lambda$。
因此
\[
b\left|\frac{\mathrm dn}{\mathrm d\lambda}\right|
\delta\lambda\simeq\lambda,
\qquad
\boxed{\frac{\lambda}{\delta\lambda}
=b\left|\frac{\mathrm dn}{\mathrm d\lambda}\right|}.
\]
\excheck 这是假设棱镜的分辨受其底边有效光束口径限制的\emph{衍射极限}；与仅给角散度的公式不同。

\diffractionproblem{144}{211}
光谱仪的三棱镜玻璃色散满足
$n=A+B/\lambda^2$，其中
$B=0.010\,\mathrm{\mu m^2}$。
(a) 求极限分辨本领随波长的关系；棱镜底边
$b=5.0$cm 时在
$\lambda_1=434$nm、
$\lambda_2=656$nm 处的数值；
(b) 分辨钠黄双线
$589.0,589.6$nm 所需的底边宽度。
\exsolution
对色散关系求导
$\mathrm dn/\mathrm d\lambda=-2B/\lambda^3$。
由第143题
\[
\boxed{\mathcal R=\frac{\lambda}{\delta\lambda}
=\frac{2bB}{\lambda^3}},
\]
说明 $\mathcal R\propto\lambda^{-3}$。
单位一致地取 $b=50000\,\mathrm{\mu m}$：
\[
\boxed{\mathcal R(0.434\,\mathrm{\mu m})
\simeq1.2\times10^4},\quad
\boxed{\mathcal R(0.656\,\mathrm{\mu m})
\simeq0.35\times10^4}.
\]
(b) 分辨条件
$2bB/\lambda^3\ge\lambda/\delta\lambda$，
所以
\[
\boxed{b_{\min}=
\frac{\lambda^4}{2B\delta\lambda}}.
\]
取 $\lambda\simeq0.589\,\mathrm{\mu m}$、
$\delta\lambda=0.0006\,\mathrm{\mu m}$，
$\boxed{b_{\min}\simeq1.0\,\mathrm{cm}}$。
\excheck $B$ 的单位为 $\mathrm{\mu m^2}$；
若把 $b$ 的厘米数直接当微米使用，会使分辨本领差四个数量级。

\diffractionproblem{145}{211}
若棱镜在目标波长附近
$|\mathrm dn/\mathrm d\lambda|
=0.10\,\mathrm{\mu m^{-1}}$，
求棱镜底边宽度，
使它具有一片刻线数
$N=10000$ 的光栅在二级谱中的相同分辨本领。
\exsolution
二级光栅分辨本领
$\mathcal R_{\rm gr}=mN=2(10000)=20000$。
棱镜衍射极限分辨本领
$\mathcal R_{\rm prism}=b|\mathrm dn/\mathrm d\lambda|$。
令二者相同：
\[
\boxed{b=
\frac{mN}{|\mathrm dn/\mathrm d\lambda|}
=\frac{20000}{0.10}\,\mathrm{\mu m}
=0.20\,\mathrm m=20\,\mathrm{cm}}.
\]
\excheck 这里采用棱镜\emph{底边}长度而非光路宽的任意投影，光栅亦以实际被照亮的刻线数计算。

\diffractionproblem{146}{211}
物镜直径 $D=5.0$cm 的望远镜，取
$\lambda=0.55\,\mathrm{\mu m}$。
求物镜的角分辨本领以及距望远镜
$\ell=3.0$km 的两物点的最小可分辨间距。
\exsolution
圆孔艾里斑瑞利判据
$\delta\psi_{\min}=1.22\lambda/D$，
角分辨本领
\[
\boxed{\mathcal R_\psi=
\frac1{\delta\psi_{\min}}
=\frac D{1.22\lambda}
\simeq7.4\times10^4}.
\]
远处两点最小间距
$\Delta y\simeq\ell\delta\psi_{\min}$，
所以
\[
\boxed{\Delta y_{\min}
\simeq\ell\frac{1.22\lambda}{D}
=\frac{3000(1.22)(0.55\times10^{-6})}{0.050}
\simeq4.0\,\mathrm{cm}}.
\]
\excheck 这里的“分辨本领”按\emph{最小角度的倒数}定义，而不是线性放大率。

\diffractionproblem{147}{211}
口径 $D=5$m 的反射望远镜在
$\lambda=0.55\,\mathrm{\mu m}$ 处工作。
求在月面上能分辨的两个点之间的最小距离。
\exsolution
由瑞利判据最小角距
$\delta\psi=1.22\lambda/D$。
取平均地月距离
$L_{\rm Moon}\simeq3.84\times10^8$m，
则
\[
\boxed{\Delta y_{\min}=
L_{\rm Moon}\delta\psi
=\frac{(3.84\times10^8)(1.22)(0.55\times10^{-6})}{5}
\simeq5.2\times10^1\,\mathrm m
\approx50\,\mathrm m}.
\]
\excheck 此为理想衍射极限；地面大气视宁度通常会劣化实际角分辨率，不能直接当成现实观测保证。

\diffractionproblem{148}{211}
望远镜物镜口径
$D=5.0$cm，观察者瞳孔直径
$d_0=4.0$mm。
求完全利用物镜理论分辨本领所需的最小角放大率。
\exsolution
物镜的瑞利最小角距
$\delta\psi=1.22\lambda/D$，
人眼瞳孔对同波长的理想衍射分辨极限
$\delta\psi_{\rm eye}=1.22\lambda/d_0$。
望远镜使角距放大 $\Gamma$，
使
$\Gamma\delta\psi\ge\delta\psi_{\rm eye}$。
因此
\[
\boxed{\Gamma_{\min}=
\frac{\delta\psi_{\rm eye}}{\delta\psi}
=\frac D{d_0}
=\frac{50\,\mathrm{mm}}{4.0\,\mathrm{mm}}
=12.5\approx13}.
\]
\excheck 这里的最小放大率是衍射分辨利用条件，
不是在所有目标亮度下的最佳视觉放大率。

\diffractionproblem{149}{212}
显微镜物镜的数值孔径
$\mathrm{NA}=\sin\alpha=0.24$，
用波长
$\lambda=0.55\,\mathrm{\mu m}$ 的光照亮物体。
求该显微镜能分辨的两物点最小间距。
\exsolution
圆孔物镜在物方的瑞利横向分辨极限
\[
\boxed{d_{\min}=
\frac{0.61\lambda}{\mathrm{NA}}
=\frac{0.61(0.55\,\mathrm{\mu m})}{0.24}
\simeq1.4\,\mathrm{\mu m}}.
\]
\excheck 原题给的数值孔径已是 $\sin\alpha$，
不可再对0.24取一次正弦；
这里默认物体位于空气中，折射率为1。

\diffractionproblem{150}{212}
显微镜物镜数值孔径为
$\sin\alpha=0.24$，观察者瞳孔
$d_0=4.0$mm。
求使物镜衍射分辨本领被眼睛完全利用的最小放大率。
\exsolution
显微镜物方线分辨极限
$d_{\min}=0.61\lambda/\sin\alpha$。
若将它放在标准明视距离
$\ell_0=25$cm 观看，
该间距所张裸眼角为
$d_{\min}/\ell_0$。
眼瞳衍射角极限约
$1.22\lambda/d_0$。
需放大
$\Gamma_{\min}d_{\min}/\ell_0
=1.22\lambda/d_0$，故
\[
\boxed{\Gamma_{\min}=
\frac{2\ell_0\sin\alpha}{d_0}
=\frac{2(0.25)(0.24)}{0.004}
=30}.
\]
\excheck 公式中波长消掉，这只是将物镜极限角尺度匹配到眼瞳的衍射极限；视网膜和视觉神经等实际因素未纳入。
